% ============================
%   3. Related work
% ============================
\section{Related work}
\label{sec:related_work}

\subsection{Production-grade NIDS}

\textbf{Snort}~\cite{roesch1999snort} is the original lightweight
signature IDS and remains the de-facto reference. Its strengths are
maturity, a vast rule ecosystem (Emerging Threats, Talos), and a
permissive license. Its main shortcoming for the use case targeted here
is that it is a console-only daemon: integrating Snort into an analyst
workflow requires either a commercial console (Snort Cloud) or a third
party such as SecurityOnion.

\textbf{Suricata}~\cite{suricata} re-engineered Snort for multi-threaded
operation and added protocol-aware parsing. It remains rule-driven and,
like Snort, ships no built-in dashboard.

\textbf{Zeek (formerly Bro)}~\cite{paxson1999bro} occupies a different
niche: it logs richly-structured events about every flow and lets
operators write detection logic in a scripting language. It is widely
used in research but, again, has no embedded analyst UI.

The common gap shared by all three is forensic context: alerts are
emitted with timestamp and 5-tuple, but the raw packets that triggered
them are typically not retained. Operators bolt on full packet capture
(tcpdump, Moloch/Arkime) to fill that gap. \nidsname{} treats forensics
as a first-class concern of the detection engine itself.

\subsection{Machine-learning IDS}

The literature on applying machine learning to intrusion detection is
vast; recent surveys list hundreds of distinct
proposals~\cite{buczak2016survey,khraisat2019survey}. Most of this work
focuses on benchmark accuracy on synthetic datasets:
KDD'99~\cite{tavallaee2009detailed}, the cleaned NSL-KDD, and more
recent dataset releases such as CICIDS2017~\cite{sharafaldin2018cicids}.

A persistent critique of this line of research is that benchmark
accuracy does not transfer to operational deployment. Sommer and
Paxson~\cite{sommer2010outside} argue that the closed-world assumption
of standard ML evaluation does not hold in networking, where the
distribution of benign traffic itself drifts continuously. They
recommend (i) anomaly detection rather than misuse detection for the ML
component, (ii) operator-friendly explanations of why a model fired, and
(iii) keeping the ML engine in an advisory role behind a signature
engine. \nidsname{} embraces all three recommendations: Isolation Forest
is offered as a first-class option, every alert is labelled with its
detection method, and hybrid agreement is highlighted prominently in the
UI.

\subsection{SIEM and forensics platforms}

\textbf{Wazuh}, \textbf{SecurityOnion} and \textbf{Splunk} are
fully-fledged SIEM platforms with analyst consoles. They are too large
to deploy on a single laptop and they treat detection as an
orchestration problem (collecting alerts from external engines) rather
than as a unified pipeline. \nidsname{} occupies the much smaller niche
of \textit{a single-binary teaching/demo SIEM with detection and
forensics built in}, intentionally trading scale for clarity.

\subsection{Positioning of this work}

\begin{table}[h]
\centering
\small
\begin{tabularx}{\textwidth}{l c c c c X}
\toprule
& Snort & Suricata & Zeek & ML-IDS research & \textbf{\nidsname{}} \\
\midrule
Signature engine            & yes & yes & partial & rare & \textbf{yes} \\
ML engine                   & no  & no  & no      & yes & \textbf{yes (RF, IF)} \\
Hybrid labelling            & --  & --  & --      & --  & \textbf{yes} \\
Embedded analyst UI         & no  & no  & no      & no  & \textbf{yes} \\
Per-alert packet capture    & no  & no  & no      & no  & \textbf{yes} \\
Single-binary deployment    & no  & no  & no      & --  & \textbf{yes (laptop demo)} \\
\bottomrule
\end{tabularx}
\caption{Positioning of \nidsname{} against existing tool categories.}
\label{tab:positioning}
\end{table}

The intention is not to outperform mature production engines on raw
throughput or rule count, but to demonstrate, in a single contained
application, that the three ingredients --- signature, ML and forensics
--- can be made to cooperate inside a tidy SIEM-style user experience.
